% \newpage
% \pagebreak

\section{Proof of Theorem 1}
\label{sec: thm1}

% \textbf{Theorem 1.} Given the input prompt $x\sim \mathcal X$, let $y_c\sim \mathcal Y_c$ and $y_p\sim \mathcal Y_p$ denotes the clean and poisoned responses to $x$, respectively. 
% $p_{\theta}(\cdot|x)$ is the output probability with an LLM parameterzed by $\theta$.
% $P_{posi}(x;\theta)$ and  $P_{clean}(x;\theta)$ are the probability of generating poison and clean outputs from $x$.
% $\mathbb H_{\theta}(\cdot|x)$ is the output entropy of $p_{\theta}(\cdot|x)$.
% % The following upperbound holds for the DPO objective \cite{rafailov2024direct} $\mathcal L_{DPO}$,
% The DPO objective \cite{rafailov2024direct} $\mathcal L_{DPO}$ can be upperbounded by $\mathcal L_{DPO}^u$, %\emph{s.t.},
% % %
% % \begin{equation}
% %     \mathcal L_{DPO} ^u =  \mathbb E_{(x,y_c,y_p)\sim \mathcal D} \left( -\frac{\mathbb H(y_p|x)}{P_{posi}(x)} + \frac{\mathbb H(y_c|x)}{P_{clean}(x)} \right) 
% % \end{equation}
% % %
% %
% \begin{equation}
% \begin{split}
%     \mathcal L_{DPO} ^u =& \\ \mathbb E_{(x,y_c,y_p)\sim \mathcal D} &\left( -\frac{\mathbb H(y_p|x)}{P_{posi}(x)} + \frac{\mathbb H(y_c|x)}{P_{clean}(x)} \right) + \mathcal C_{ref, \mathcal D}
% \end{split}
% \end{equation}
% %
% with the following assumptions on $(x, y_c, y_p)\sim \mathcal D = (X, Y_c, Y_p)$,
% %
% \begin{align}
%     y^* = \text{argmax}&_y p(y|x)\,\, \text{is poisoned.} \\
%     p(y_p|x) >& \,\,p(y_c|x)
% \end{align}
% %






% \begin{table}[t!]
% % \begin{wraptable}{r}{0.6\textwidth}
% \centering
% \small

%  \begin{minipage}[b]{\linewidth}
% \resizebox{\linewidth}{!}{
% \begin{tabular}{l | c | c c c}
% \toprule
%  & \textbf{w/ $\Phi$} & \textbf{ASR} \downarrow & \textbf{F1 (clean)} \uparrow & \textbf{F1 (attack)} \uparrow \\
% \midrule
% \multirow{2}{*}{$p$=0.5\%}    & Y           & \textbf{7.44} $\pm$ 0.22  & 28.68 $\pm$ 0.30  & \textbf{22.84} $\pm$ 0.49 \\
%   & N    & 7.65 $\pm$ 0.63 & \textbf{29.06} $\pm$0.42 & 21.35 $\pm$ 0.44 \\
%   \midrule
% \multirow{2}{*}{$p$=1.0\%}    & Y           & 4.38 $\pm$ 0.32  & \textbf{28.12} $\pm$ 0.58  & \textbf{23.17} $\pm$ 0.30 \\
%   & N    & \textbf{4.64} $\pm$ 0.53 & 26.01 $\pm$ 0.43 & 18.77 $\pm$ 0.91 \\
%   \midrule

%   \multirow{2}{*}{$p$=5.0\%}    & Y           & \textbf{0.83} $\pm$ 0.07  & \textbf{27.48} $\pm$ 0.66 & \textbf{23.03} $\pm$ 0.64\\
%   & N    & 7.86 $\pm$ 0.86 & 6.48 $\pm$ 1.04 & 10.43 $\pm$ 2.05 \\
%   \midrule
% \bottomrule
% \end{tabular}
% }
% % \vspace{-1em}
% \caption{
% Ablation on whether to prune from $\Phi$ using SQuAD with LLaMA-3-8B.
% “Y” means that \eqref{eq: tao} and \eqref{eq: mask} only mask and prune the neurons defined in \eqref{eq: phi}. 
% “N” means that $\Phi$ represents all neurons in the KV cache.
% Pruning from $\Phi$ becomes increasingly important as more neurons are pruned ($p \uparrow$), \emph{i.e.}, when robustness requirements become stricter relative to utility.
% }
% \label{tb:phi}
% \end{minipage}

% \vspace{-1em}
% \end{table}

\begin{table}[t!]
\centering
\small
\begin{minipage}[b]{\linewidth}
\resizebox{\linewidth}{!}{
\begin{tabular}{l | c | c c c}
\toprule
 & \textbf{w/ $\Phi$} & \textbf{ASR}~$\downarrow$ & \textbf{F1 (clean)}~$\uparrow$ & \textbf{F1 (attack)}~$\uparrow$ \\
\midrule
\multirow{2}{*}{$p$=0.5\%}    & Y           & \textbf{7.44} $\pm$ 0.22  & 28.68 $\pm$ 0.30  & \textbf{22.84} $\pm$ 0.49 \\
  & N    & 7.65 $\pm$ 0.63 & \textbf{29.06} $\pm$0.42 & 21.35 $\pm$ 0.44 \\
\midrule
\multirow{2}{*}{$p$=1.0\%}    & Y           & 4.38 $\pm$ 0.32  & \textbf{28.12} $\pm$ 0.58  & \textbf{23.17} $\pm$ 0.30 \\
  & N    & \textbf{4.64} $\pm$ 0.53 & 26.01 $\pm$ 0.43 & 18.77 $\pm$ 0.91 \\
\midrule
\multirow{2}{*}{$p$=5.0\%}    & Y           & \textbf{0.83} $\pm$ 0.07  & \textbf{27.48} $\pm$ 0.66 & \textbf{23.03} $\pm$ 0.64\\
  & N    & 7.86 $\pm$ 0.86 & 6.48 $\pm$ 1.04 & 10.43 $\pm$ 2.05 \\
\midrule
\bottomrule
\end{tabular}
}
\caption{
Ablation on whether to prune from $\Phi$ using SQuAD with LLaMA-3-8B.
“Y” means that \eqref{eq: tao} and \eqref{eq: mask} only mask and prune the neurons defined in \eqref{eq: phi}. 
“N” means that $\Phi$ represents all neurons in the KV cache.
Pruning from $\Phi$ becomes increasingly important as more neurons are pruned ($p \uparrow$), \emph{i.e.}, when robustness requirements become stricter relative to utility.
}
\label{tb:phi}
\end{minipage}
\vspace{-1em}
\end{table}





% \begin{table}[t!]
% % \begin{wraptable}{r}{0.6\textwidth}
% \centering
% \small

%  \begin{minipage}[b]{\linewidth}
% \resizebox{\linewidth}{!}{
% \begin{tabular}{l | c | c c c}
% \toprule
%  & \textbf{$N$} & \textbf{ASR} \downarrow & \textbf{F1 (clean)} \uparrow & \textbf{F1 (attack)} \uparrow \\
%  \midrule
%  \multirow{3}{*}{Llama3-8B}    & $4$   &     8.15 $\pm$ 0.58 & 26.14 $\pm$ 0.76 & 22.99 $\pm$ 1.19 \\
%   & $8$                                    & 7.44 $\pm$ 0.22 & 28.68 $\pm$ 0.30 & 22.84 $\pm$ 0.49\\
%    & $12$  &                                 7.52 $\pm$ 0.31 & 28.27 $\pm$ 0.42 & 24.12 $\pm$ 0.50	 \\
% \midrule
% \multirow{3}{*}{Mistral-7B}    & $4$          & 0.74 $\pm$ 0.36 & 24.75 $\pm$ 0.65 & 22.52 $\pm$ 1.67 \\
%   & $8$     & 0.68 $\pm$ 0.41 & 24.46 $\pm$ 0.91 & 23.10 $\pm$ 1.32\\
%    & $12$   & 0.57 $\pm$ 0.32 & 25.05 $\pm$ 1.05 & 23.26 $\pm$ 0.88	 \\
%    \midrule
% \multirow{3}{*}{\shortstack{Phi-3.5-mini-\\instruct (3.8B)}}    & $4$           & 0.86 $\pm$ 0.32 & 25.76 $\pm$ 0.77 & 27.12 $\pm$ 1.31 \\
%   & $8$     & 0.71 $\pm$ 0.18 & 26.76 $\pm$ 0.56 & 25.55 $\pm$ 0.60\\
%    & $12$   & 0.62$\pm$ 0.13 & 26.37 $\pm$ 0.71 & 25.26 $\pm$ 0.47	 \\

% \bottomrule
% \end{tabular}
% }
% % \vspace{-1em}
% \caption{
% Performance with number of samples $N$ used for attribution. We observe that the ASR generally decreases as 
% N increases, with the variance also showing decreasing trend. 
% % Performance of LLama3-8b on SQuAD with differen values of $\alpha$. 
% % $\mathcal L^{attr}_{full}$ means we attribute with all the tokens in the response.
% }
% \label{tb:N}
% \end{minipage}

% \vspace{-1em}
% \end{table}



\begin{table}[t!]
\centering
\small

\begin{minipage}[b]{\linewidth}
\resizebox{\linewidth}{!}{
\begin{tabular}{l | c | c c c}
\toprule
 & \textbf{$N$} & \textbf{ASR}~$\downarrow$ & \textbf{F1 (clean)}~$\uparrow$ & \textbf{F1 (attack)}~$\uparrow$ \\
\midrule
\multirow{3}{*}{Llama3-8B}    
  & $4$  & 8.15 $\pm$ 0.58 & 26.14 $\pm$ 0.76 & 22.99 $\pm$ 1.19 \\
  & $8$  & 7.44 $\pm$ 0.22 & 28.68 $\pm$ 0.30 & 22.84 $\pm$ 0.49\\
  & $12$ & 7.52 $\pm$ 0.31 & 28.27 $\pm$ 0.42 & 24.12 $\pm$ 0.50 \\
\midrule
\multirow{3}{*}{Mistral-7B}    
  & $4$  & 0.74 $\pm$ 0.36 & 24.75 $\pm$ 0.65 & 22.52 $\pm$ 1.67 \\
  & $8$  & 0.68 $\pm$ 0.41 & 24.46 $\pm$ 0.91 & 23.10 $\pm$ 1.32\\
  & $12$ & 0.57 $\pm$ 0.32 & 25.05 $\pm$ 1.05 & 23.26 $\pm$ 0.88 \\
\midrule
\multirow{3}{*}{\shortstack{Phi-3.5-mini-\\instruct (3.8B)}}    
  & $4$  & 0.86 $\pm$ 0.32 & 25.76 $\pm$ 0.77 & 27.12 $\pm$ 1.31 \\
  & $8$  & 0.71 $\pm$ 0.18 & 26.76 $\pm$ 0.56 & 25.55 $\pm$ 0.60\\
  & $12$ & 0.62 $\pm$ 0.13 & 26.37 $\pm$ 0.71 & 25.26 $\pm$ 0.47 \\
\bottomrule
\end{tabular}
}
\caption{
Performance with number of samples $N$ used for attribution. 
We observe that the ASR generally decreases as $N$ increases, 
with the variance also showing a decreasing trend.
}
\label{tb:N}
\end{minipage}

\vspace{-1em}
\end{table}







\noindent\textbf{Theorem 1.} Given the input prompt $x\sim \mathcal X$, let $y^c\sim \mathcal Y^c_x$ and $y^p\sim \mathcal Y^p_x$ denotes the clean and poisoned responses to $x$, respectively. 
$(x, y^c, y^p)\sim \mathcal D = (X, Y^c_x, Y^p_x)$ is the dataset of perference optimization.
$p_{\theta}(\cdot|x)$ is the output probability with an LLM parameterzed by $\theta$.
% If we hold the following assumptions on $(x, y_c, y_p)\sim \mathcal D = (X, Y_c, Y_p)$,
% %
% \begin{align}
%     y^* = \text{argmax}&_y p(y|x)\,\, \text{is poisoned.} \label{eq:a1}\\
%     p(y_p|x) >& \,\,p(y_c|x) \label{eq:a2}
% \end{align}
% %
The preference optimization with $ \mathcal L_{DPO}$ can be upperbounded by $\mathcal L_{DPO}^u$, \emph{s.t.},
%
\begin{equation}
  \begin{split} \label{eq: L_dpo_u_p}
    \mathcal L^u_{DPO} &= \mathbb E_{x\sim \mathcal X } ( \,\,\log \frac{p_\theta(y\in |\mathcal Y^p_x|\,|x)}{p_\theta(y\in |\mathcal Y^c_x|\,|x)} \\
    &+ \mathbb H(\mathcal Y^c_x|x) - \mathbb H(\mathcal Y^p_x|x) \,\,) + \mathcal C_{ref, \mathcal D}
  \end{split}
\end{equation}
%
where $|\cdot|$ is the support of a distribution. $\mathcal C_{ref, \mathcal D}$ is a constant to $\theta$ that is functioned by $\mathcal D$ and the DPO reference model $ref$.
% $P_{posi}(x;\theta)$ and  $P_{clean}(x;\theta)$ are the probability of generating poison and clean outputs from $x$.
$\mathbb H(\mathcal Y^c_x|x)$ and  $\mathbb H(\mathcal Y^p_x|x)$, respectively, are the entropy of clean and poisoned responses given $x$.
$p_\theta(y\in |\mathcal Y^p_x|\,|x)$ is the probability of generating poisoned responses from $x$, and similar to $p_\theta(y\in |\mathcal Y^c_x|\,|x)$. 

\vspace{2mm}
\noindent\textbf{Proof.}
In the context of defending against the prompt injection attack with $(x,y^c,y^p)\sim \mathcal D$, the DPO objective $\mathcal L_{DPO}$ can be defined as,
%
\begin{equation}
  \begin{split} \label{eq:L_dpo}
    \mathcal L_{DPO} = \mathbb  E&_{(x,y^c,y^p)\sim \mathcal D} [ \log \,\sigma( \beta\log \frac{p_{\theta}(y^p|x)}{p_{ref}(y^p|x)} \\
    &-\beta\log \frac{p_{\theta}(y^c|x)}{p_{ref}(y^c|x)})].
  \end{split}
\end{equation}
%
where $\sigma(\cdot)$ is the sigmoid function and $\beta>0$ is a regularization parameter.
The reference model $ref$ serves as an anchor in a way that the minimization of $\mathcal L_{DPO}$ is also minimizing the following KL divergence,
%
\begin{equation} \label{eq:DPO_KL}
    \mathcal D_{KL} [p_{\theta}(y|x)\,||\,p_{ref}(y|x)].
\end{equation}
%
$ref$ is chosen before training. In this proof, we choose $ref$ to be a model that is more immune to the indirect prompt injection attack compared to the LLM with $\theta$, \emph{i.e.},
%
\begin{align} % \label{eq:a_ref}
    p_{ref}(y^c|x) > p_{\theta}(y^c|x)  \label{eq:a_ref_1}\\
    p_{ref}(y^p|x) > p_{\theta}(y^c|x)   \label{eq:a_ref_2}
\end{align}
%
This choice of $ref$ is reasonable since it makes $\mathcal L_{DPO}$ a strong object in defending against prompt injection attack due to \eqref{eq:DPO_KL}.

\vspace{2mm}
\noindent With \eqref{eq:a_ref_1} and \eqref{eq:a_ref_1}, we can observe that,
%
\begin{equation} \label{eq:sigma_inputs}
    \mathcal S  = \log \frac{p_{\theta}(y_p|x)}{p_{ref}(y_p|x)} 
    -\log \frac{p_{\theta}(y_c|x)}{p_{ref}(y_c|x)}) > 0
\end{equation}
%
This follows that the $\log\sigma(\cdot)$ in \eqref{eq:L_dpo} should be concave since,
\begin{itemize}
    \item $\log(\cdot)$ is a concave function, and the $\sigma(\cdot)$  in \eqref{eq:L_dpo} is also concave given that its inputs $\mathcal S$ and $\beta$ are both positive.
    % \item $\log(\cdot)$ is also convex.
    \item Both  $\log(\cdot)$ and  $\sigma(\cdot)$ are monotonically increasing.
\end{itemize}



% This implies that the $\sigma(\cdot)$ in \eqref{eq:L_dpo} is a convex function since its inputs are positive given \eqref{eq:sigma_inputs} and $\beta>0$. Considering that the $log(\cdot)$ is also convex and both $log(\cdot)$ and $\sigma(\cdot)$ are monotonically increasing, their composition $\log\sigma(\cdot)$ should also be convex in \eqref{eq:L_dpo}.
\noindent Then, we can upperbound $\mathcal L_{DPO}$ following the Jensen's Inequality,
%
% \begin{equation}
\begin{align}
  \mathcal L_{DPO}  &= \mathbb  E_{(x,y^c,y^p)\sim \mathcal D} [ \log \,\sigma( \beta \cdot \mathcal S)] \\
  &\leq \log\sigma(\beta \cdot E_{(x,y^c,y^p)\sim \mathcal D} \mathcal S). \label{eq:L_u_DPO}
\end{align}
% \end{equation}
%
Since \eqref{eq:L_u_DPO} only relies on the expectation term within $\sigma(\cdot)$, we define our upperbound objective as, %\emph{i.e.},
%
\begin{align}
    \mathcal L^u_{DPO} := \mathbb E_{(x,y_c,y_p)\sim \mathcal D} \mathcal \,S 
\end{align}
%
%
% \begin{equation}
% \begin{align}
%     \mathcal L^u_{DPO} &:= E_{(x,y_c,y_p)\sim \mathcal D} \mathcal S \\
%     &= E_{(x,y_c,y_p)\sim \mathcal D} \log \frac{p_{\theta}(y_p|x)}{p_{ref}(y_p|x)} 
%     -\log \frac{p_{\theta}(y_c|x)}{p_{ref}(y_c|x)}) \\
%     =
% \end{align}
% \end{equation}
%

\noindent We rewrite $ \mathcal L^u_{DPO}$ as,
%
% \begin{equation}
% \begin{split}
\begin{align}
         \mathcal L&^u_{DPO} = \mathbb E_{(x,y^c,y^p)\sim \mathcal D} ( \,\log \frac{p_{\theta}(y^p|x)}{p_{ref}(y^p|x)} \notag\\
    &\,\,\,\,\,\,\,\,\,\,\,\,\,\,\,\,\,\,\,\,\,\,\,\,\,\,\,\,\,\,\,\,\,\,\,\,\,\,\,\,\,\,\,\,\,\,\,\,\,\,\,\,-\log \frac{p_{\theta}(y^c|x)}{p_{ref}(y^c|x)}) \\
    &=\mathbb E_{(x,y^c,y^p)\sim \mathcal D} (\log p_{\theta}(y^p|x) \notag\\
    &\,\,\,\,\,\,\,\,\,\,\,\,\,\,\,\,\,- \log p_{\theta}(y^c|x)) + \mathcal C_{ref,\mathcal D}, \label{eq:L_u_DPO_1}
\end{align}

\noindent where,
%
\begin{equation}
    \mathcal C_{ref,\mathcal D} = \mathbb E_{(x,y^c,y^p)\sim \mathcal D} \log \frac{p_{ref}(y^c|x)}{p_{ref}(y^p|x)},
\end{equation}
%
is a constant to $\theta$ that only depends on dataset $\mathcal D$ and the choice of $ref$.

\vspace{2mm}
\noindent The first term in \eqref{eq:L_u_DPO_1} can be decomposed by,
%
\begin{align}
    &\mathbb E_{(x,y^c,y^p)\sim \mathcal D} (\log p_{\theta}(y^p|x) - \log p_{\theta}(y^c|x)) \notag \\
    &\!\!\!= \mathbb E_{x\sim\mathcal X} ( \,\,\underbrace{\sum_{y^p}  p_\theta(\mathcal Y^p_x=y^p|x) \log p_{\theta}(y^p|x)}_{\mathcal V^p} \notag \\
    &- \underbrace{\sum_{y^c}   p_\theta(\mathcal Y^c_x=y^c|x)  \log p_{\theta}(y^c|x)}_{\mathcal V^c} ). \label{eq:L_u_DPO_2}
\end{align}
%
% \begin{align}
%     &\mathbb E_{(x,y_c,y_p)\sim \mathcal D} (\log p_{\theta}(y_p|x) - \log p_{\theta}(y_c|x)) \notag \\
%     &\!\!\!= \mathbb E_{x\sim\mathcal X} ( \underbrace{\sum_{y_p}  (\sum_{y_c} p(\mathcal Y_p=y_p,\mathcal Y_c=y_c|x) \log p_{\theta}(y_p|x)}_{\mathcal V_p} \notag \\
%     &- \underbrace{\sum_{y_c}  (\sum_{y_p} p(\mathcal 
%  Y_p=y_p,\mathcal  Y_c=y_c|x) )  \log p_{\theta}(y_c|x)}_{\mathcal V_c} ).
% \end{align}

% \noindent With the assumptions of \eqref{eq:a1} and \eqref{eq:a2}, it should be noted that $Y_p$ and $Y_c$ are \textbf{not} independent with given $x\in \mathcal X$. 
% Therefore, we cannot simplify \eqref{eq:L_u_DPO_2} by independently assigning values to the marginal probabilities of $p(\mathcal Y_p=y_p|x)$ and $p(\mathcal Y_c=y_c|x)$.
% Consequently, $\mathcal V_p$ and $\mathcal V_c$ in \eqref{eq:L_u_DPO_2} depends on our definition of $p(\mathcal Y_p, \mathcal Y_c | x)$  \emph{i.e.}, the sampling strategy of $(\mathcal Y_p, \mathcal Y_c | x)$ under the assumptions of \eqref{eq:a1} and \eqref{eq:a2}. 

% \vspace{2mm}
% \noindent In \textbf{Lemma 1}, we show that there exists a sampling strategy of $(\mathcal Y_p, \mathcal Y_c | x)$ that satisfies \eqref{eq:a1} and \eqref{eq:a2}, while rendering the marginal distributions, \emph{s.t.},
% %
% \begin{align}
%   \!\!\!\!\!\!\!\!\!\!p(\mathcal Y_p=y_p|x) = p_{\theta}(\mathcal Y_p=y_p|x), \\
%    \!\!\!\!\!\!\!\!\!\!p(\mathcal Y_c=y_c|x) = p_{\theta}(\mathcal Y_c=y_c|x).
% \end{align}
% %

\noindent Then, we can have,
%
\begin{align}
    \mathcal V^p &= \sum_{y^p} p_{\theta}(\mathcal Y^p_x=y^p|x) \log p_{\theta}(y^p|x) \\
    &= \sum_{y^p_x} p_{\theta}(\mathcal Y^p_x=y^p|x)\cdot\notag \\ &\!\,\,\,\,\,\,\,\,\,\,\log (p_{\theta}(\mathcal Y^p_x=y_p|x)\cdot p(y\in|\mathcal Y^p|\,|x)) \\
    % &=\sum_{y_p} p_{\theta}(\mathcal Y_p=y_p|x)\log (p_{\theta}(\mathcal Y_p=y_p|x) +\sum_{y_p} p_{\theta}(\mathcal Y_p=y_p|x) \log p(y\in|\mathcal Y_p|\,|x) \\
    &= -\mathbb H(\mathcal Y^p|x) + \log p(y\in|\mathcal Y^p|\,|x) \label{eq:V_p}
\end{align}
%
% \noindent where $P_{posi(x)}$ is the probability that the model generates poisoned response given $x$,
% %
% \begin{equation}
%     \!\!\!P_{posi}(x) = p(y\in |\mathcal Y_p| \,|x) =\sum_{y\in |\mathcal Y_p|} \, p_{\theta}(y|x)
% \end{equation}
% %
% \noindent $|\cdot|$ is the support of a distribution. In this way, $\mathcal V_p$ can be expresses as, 
% %
% \begin{align}
%     \mathcal V_p = 
% \end{align}
% %

% \end{split}
% \end{equation}
    
\noindent Similarly,  $\mathcal V^c$ can be expressed as,
%
\begin{align}
    \mathcal V^c = -\mathbb H(\mathcal Y^c_x|x) + \log p(y\in|\mathcal Y^c_x|\,|x) \label{eq:V_c}
\end{align}
%
Combining \eqref{eq:L_u_DPO_1}, \eqref{eq:V_p} and \eqref{eq:V_c} together, we can write $\mathcal L_{DPO}^u$ as,
%
\begin{equation}
  \begin{split}
    \mathcal L^u_{DPO} &= \mathbb E_{x\sim \mathcal X } ( \,\,\log \frac{p_\theta(y\in |\mathcal Y^p_x|\,|x)}{p_\theta(y\in |\mathcal Y^c_x|\,|x)} \\
    &+ \mathbb H(\mathcal Y^c_x|x) - \mathbb H(\mathcal Y^p_x|x) \,\,) + \mathcal C_{ref, \mathcal D}
  \end{split}
\end{equation}
%
